\begin{proposition}[The first missing frequency]\label{prop:firstres}
The least positive frequency missed by any four-point shape is the number
$\omega_*$ enclosed in Section~\ref{sec:class}. Thus a convergent four-point
variance restricts the half-width of the limiting normalized zeta counts to
\[
 \frac{0.1604}{\sqrt{1+\omega_*^2}}=0.025405106245\ldots .
\]
This constant is sharp for spectra satisfying the two local bounds used in~\eqref{eq:cap}.
\end{proposition}
\begin{proof}
In the triangle interval, $A=B/2+C/3$, with $B,C>0$. A zero therefore requires
$2^{\ii\omega}-1/2$ and $3^{\ii\omega}-1/3$ to point in opposite directions.
Write $L=2\pi/\log3$ and $U=3\pi/(2\log2)$. For $0<\omega\le L$ the arguments
$\omega\log2$ and $\omega\log3$ bound an arc of length at most
$2\pi\log(3/2)/\log3<\pi$. The positive combination
$B2^{\ii\omega}+C3^{\ii\omega}$ points strictly inside that arc, whose interior excludes
the positive real axis. Hence this range contains no missing frequency.

Let $\theta_p(\omega)$ be the continuous argument of
$p^{\ii\omega}-1/p$ starting at zero. With $\phi=\omega\log p$,
\[
 \theta'_p=(\log p)\frac{1-p^{-1}\cos\phi}{1-2p^{-1}\cos\phi+p^{-2}}.
\]
On $[L,U]$ the angle $\omega\log2$ lies between $\pi$ and $3\pi/2$,
so $\theta'_2\le\log2$. The same formula gives
$\theta'_3\ge(3/4)\log3>\log2$. Thus $\theta_3-\theta_2$ strictly increases.
At $L$ its value lies in $(0,\pi)$ since $\theta_3(L)=2\pi$ and
$\pi<\theta_2(L)<2\pi$. At $U$ we have $\theta_2(U)=\pi+\arctan 2$,
whereas $\theta_3(U)=2\pi+\arg(\ee^{\ii\psi}-1/3)$ with
$\psi=(3\alpha/2-2)\pi\in(\pi/3,\pi/2)$, using $3^9>2^{14}$
and $3^3<2^5$. Here $\sin\psi>2(\cos\psi-1/3)$, so
$\pi<\theta_3(U)-\theta_2(U)<2\pi$. The unique crossing of $\pi$
is the first possible opposition. The enclosed pair $(t_*,\omega_*)$
already realizes a zero in $(L,U)$, hence gives this crossing.
Equation~\eqref{eq:cap} gives $r\le(1.3208-1)/2$;
the discrete construction in Section~\ref{sec:discrete} attains equality at $\omega_*$.
\end{proof}
